% Source fingerprint: 7d0f5040f7fc6d6fea8100161aa640173e918d11b9239f8fbb3802626d2eee38
% T03: raw TeX cells preserved; no numeric highlighting.
\begin{table}[htbp]
\centering\small\renewcommand{\arraystretch}{1.18}\setlength{\tabcolsep}{4pt}
\caption[Lusin 与 Egoroff 的量词对照]{Lusin 与 Egoroff 的量词对照。两者都允许舍弃小集，但改善的对象与性质不同；表中采用本章已证明的版本。}
\label{b1:tab:visual-t03}
\begin{tabularx}{\linewidth}{@{}>{\raggedright\arraybackslash}p{3.1cm}>{\raggedright\arraybackslash}p{4.7cm}>{\raggedright\arraybackslash}X@{}}
\toprule
比较项目 & Lusin 定理 & Egoroff 定理 \\
\midrule
输入对象 & 单个有限值可测函数 & 一列几乎处处收敛的可测函数 \\
本章背景 & 欧氏空间上的 Lebesgue 测度与正则逼近 & 有限测度空间上的收敛 \\
保留的集合 & 有限测度可测集内，删去任意小测度后取紧集 & 删去任意小测度的可测例外集 \\
得到的性质 & 原函数在保留紧集上的限制连续 & 整列在保留集上一致收敛 \\
量词重点 & 每个容许误差对应一个好的紧集 & 固定好集后，收敛指标对该集中的点统一 \\
不能直接推出 & 函数在全域连续 & 原集合上一致收敛，或无限测度下的同样结论 \\
\bottomrule
\end{tabularx}
\end{table}
